\subsection{Groupoid Perspective: Orbits of Value-Preserving Rewriting and the Zeckendorf Transversal}\label{sec:groupoid-zeckendorf}

To give an auditable formulation of ``why folding is emergence rather than injection,'' we realize the Fibonacci relations as a family of \textbf{invertible local transformations} (groupoid generators) and interpret the Zeckendorf-legal sequences as a transversal of the resulting orbits.

\begin{definition}[Fibonacci local transformations (invertible rewriting)]\label{def:fib-moves}
Define the following partially invertible transformations acting on $\mathcal{D}$ (on their respective domains of applicability):
\begin{itemize}
  \item \textbf{Baseline transformation:}
  \[
  E_0(a)=a+2e_1-e_2\quad\text{when }a_{2}\ge 1,
  \qquad
  E_0^{-1}(a)=a-2e_1+e_2\quad\text{when }a_{1}\ge 2.
  \]
  \item \textbf{Recursive transformation ($k\ge 1$):}
  \[
  E_k(a)=a+e_k+e_{k+1}-e_{k+2}\quad\text{when }a_{k+2}\ge 1,
  \]
  and its inverse $E_k^{-1}$, defined when $a_k,a_{k+1}\ge 1$, is
  \[
  E_k^{-1}(a)=a-e_k-e_{k+1}+e_{k+2}.
  \]
\end{itemize}
Let $\mathcal{G}$ be the groupoid generated by $\{E_k,E_k^{-1}\}_{k\ge 0}$ (with objects $\mathcal{D}$ and morphisms given by compositions of these partial bijections).
\end{definition}

\begin{proposition}[Orbit invariant]\label{prop:orbits-preserve-val}
For any groupoid morphism (an invertible transformation obtained by composing $E_k$ and $E_k^{-1}$) $a\mapsto a'$, we have $\mathrm{Val}(a)=\mathrm{Val}(a')$.
\end{proposition}

\textbf{Proof:}
For the generator $E_k$,
\[
\mathrm{Val}(E_k(a))-\mathrm{Val}(a)
=F_{k+1}+F_{k+2}-F_{k+3}=0,
\]
and likewise for the inverse. Compositions preserve the invariant.
For the baseline generator $E_0$, we also have $\mathrm{Val}(E_0(a))-\mathrm{Val}(a)=2F_2-F_3=0$, so the conclusion holds.

\begin{theorem}[Orbit transversal and the group-theoretic interpretation of folding]\label{thm:zeckendorf-transversal}
For any $a\in\mathcal{D}$, its $\mathcal{G}$-orbit coincides with its value fiber:
\[
\mathcal{G}\cdot a \ =\ \{b\in\mathcal{D}:\mathrm{Val}(b)=\mathrm{Val}(a)\}.
\]
Moreover, this orbit intersects $X_\infty$ in exactly one point, namely $Z(\mathrm{Val}(a))$. Therefore, defining
\[
\Fold_\infty(a):=Z(\mathrm{Val}(a))\in X_\infty
\]
yields the fold operator as ``projection along orbits onto the transversal''; it compresses each unobservable value-preserving symmetry orbit into a unique stable representative.
\end{theorem}

\textbf{Proof:}
By Proposition~\ref{prop:orbits-preserve-val}, $\mathcal{G}\cdot a$ is contained in the value fiber. The reverse inclusion follows from Theorem~\ref{thm:monoid-quotient-is-N}: if $\mathrm{Val}(b)=\mathrm{Val}(a)$, then $[b]=[a]$, i.e., $b\equiv a$. The congruence $\equiv$ is precisely the closure, on the configuration space, of the baseline congruence $e_2\equiv 2e_1$ and the recursive congruence $e_{k+2}\equiv e_{k+1}+e_k$; the corresponding invertible local rewriting rules are realized by $E_0$ and $E_k$ ($k\ge 1$), so there exists a groupoid morphism sending $a$ to $b$. Finally, the element in the value fiber that lies in $X_\infty$ is unique by Zeckendorf uniqueness, and by definition $\mathrm{Val}(Z(\mathrm{Val}(a)))=\mathrm{Val}(a)$, so $Z(\mathrm{Val}(a))$ belongs to that value fiber; hence the intersection point is unique and as stated.

\begin{remark}[Two free commutative monoid coordinate systems and the canonical section]\label{rem:two-coordinate-monoids}
The configuration space $\mathcal{D}$ is the free commutative monoid generated by $\{e_k\}_{k\ge 1}$. The value map $\mathrm{Val}$ takes the quotient by the Fibonacci relations and is isomorphic to $(\NN,+)$ (Theorem~\ref{thm:monoid-quotient-is-N}); the fold $\Fold_\infty$ provides a \textbf{canonical section} of this quotient map: each congruence class has a unique representative in $X_\infty$.
On the other hand, the multiplicative monoid $(\NN_{>0},\times)$ is the free commutative monoid generated by the set of primes: the integer $n$ corresponds to the exponent vector $(v_p(n))_p$. Thus the same ``integer object'' possesses two free coordinate systems: the Fibonacci/Zeckendorf coordinates are adapted to additive normalization, while the prime exponent coordinates are adapted to the linearization of multiplication. The coordinate transformation between them carries global arithmetic information; within the scan--projection readout framework, such nonlocal information must be borne by defect dimensions such as ``carries/collisions/multiplicities.''
\end{remark}

\subsubsection{Directed normalization rewriting (compiled formulation of \texorpdfstring{$\Fold_\infty$}{Fold\_infty})}\label{subsec:directed-rewrite}
Theorem~\ref{thm:zeckendorf-transversal} has already interpreted folding as ``projection along value-preserving orbits onto the unique stable representative.'' To obtain a directly executable normalization procedure, one directs the groupoid generators into a set of \textbf{value-preserving rewriting rules}:
\begin{itemize}
  \item \textbf{Adjacent merge rule (\ref{rule:fold:adj}):} If $a_k\ge 1$ and $a_{k+1}\ge 1$, then
  \[
  a\ \to\ a-e_k-e_{k+1}+e_{k+2}.
  \]
  \item \textbf{Base deduplication rule (\ref{rule:fold:dedup1}):} If $a_1\ge 2$, then
  \[
  a\ \to\ a-2e_1+e_2.
  \]
  \item \textbf{Second-position deduplication rule (\ref{rule:fold:dedup2}):} If $a_2\ge 2$, then
  \[
  a\ \to\ a-2e_2+e_1+e_3.
  \]
  \item \textbf{General-position deduplication rule (\ref{rule:fold:dedupk}) ($k\ge 3$):} If $a_k\ge 2$, then
  \[
  a\ \to\ a-2e_k+e_{k-2}+e_{k+1}.
  \]
\end{itemize}
These correspond respectively to
\(
F_{k+1}+F_{k+2}=F_{k+3},\
2F_2=F_3,\
2F_3=F_2+F_4,\
2F_{k+1}=F_{k-1}+F_{k+2}
\),
so each rule preserves $\mathrm{Val}$.

\begin{proposition}[Termination and unique normal form]\label{prop:directed-rewrite-termination}
The above rewriting system strongly terminates for any $a\in\mathcal{D}$, and the irreducible terms are exactly $X_\infty$; consequently, every reduction sequence terminates at a unique normal form, which equals $\Fold_\infty(a)$.
\end{proposition}

\textbf{Proof:}
Let
\[
A(a)=\sum_{k\ge 1} a_k,\qquad
B(a)=\sum_{k\ge 1} k\,a_k,\qquad
C(a)=a_2,
\qquad
\mu(a)=(A(a),B(a),C(a))
\]
and compare lexicographically. Rules (\ref{rule:fold:adj}) and (\ref{rule:fold:dedup1}) decrease $A$ by $1$; rule (\ref{rule:fold:dedupk}) ($k\ge 3$) decreases $B$ by $1$; rule (\ref{rule:fold:dedup2}) preserves $A$ and $B$ but decreases $C$ by $2$. Thus every rewriting step strictly decreases $\mu$, so the system terminates. An irreducible term must satisfy $a_k\le 1$ with no adjacent pair of $1$'s, i.e., it belongs to $X_\infty$; conversely, when $a\in X_\infty$, none of the four rules (\ref{rule:fold:adj})--(\ref{rule:fold:dedupk}) can fire. Since every rewriting step preserves $\mathrm{Val}$, the terminal form has the same value as the initial configuration; by the transversal uniqueness of Theorem~\ref{thm:zeckendorf-transversal}, the normal form is unique and equals $\Fold_\infty(a)$. \qed

\begin{remark}[Two natural targets for the residual]\label{rem:residual}
The fold $\Fold_\infty$ is the projection from a value fiber onto the transversal; hence the ``residual'' can be viewed as the groupoid morphism (rewriting path/orbit coordinate) from $a$ to $\Fold_\infty(a)$. At finite resolution $m$, this corresponds to the fold multiplicity $d_m$ in Appendix~D: a large number of microstates are compressed to the same stable type, and the lost information is precisely the distinguishing labels of these preimages.
\end{remark}

Numerical verification and statistical reproduction are provided by the script \texttt{scripts/exp\_fold\_infty\_rewrite.py}, whose output is
\texttt{artifacts/export/fold\_infty\_rewrite\_stats.json}.
